\section{Coherence and Persistence of the Selected Pairing}
\label{sec:phil-selected-coherence-proof}

\propositionheading{Coherence and persistence after pairing selection}{prop:phil-selected-pair-persistence}
In the four-cycle of \cref{conj:phil-contrast-formation}, let
\(0<\mathrm T\le1/100\), let \(M\) be one matching, and suppose
\([t_0,t_0+\mathsf h]\subseteq[t_{\mathrm L},2t_{\mathrm L}]\), with
\(\mathsf h>0\). Write
\(\nu_{\mathrm T}^M=\mathcal N(0,\mathrm T\mathsf A_M^{-1})\otimes\delta_M\).
Assume the learned and frozen recognition laws satisfy
\[
 \sup_{t_0\le t\le t_0+\mathsf h}
 \operatorname{TV}(q_t^{\mathrm L},\nu_{\mathrm T}^M)\le\delta,
 \qquad
 \operatorname{TV}(q_{t_0}^{0},\mu_{\mathrm T,1/2})\le\delta.
\]
Let \(\mathsf O_{\mathrm T,\mathsf h}^M\) be the stationary
Ornstein--Uhlenbeck path law with drift \(-\mathsf A_M\mathbf x\),
diffusion covariance \(2\mathrm TI\), and constant assignment \(M\).
For the configuration and assignment path laws on this interval,
\begin{align}
 \operatorname{TV}(\mathsf P^{\mathrm L}_{t_0,\mathsf h},
                    \mathsf O_{\mathrm T,\mathsf h}^M)
 &\le(\mathsf h+2)\delta,
 \label{eq:phil-selected-path-coupling}\\
 \operatorname{TV}\left(\mathsf P^{0}_{t_0,\mathsf h},
   \tfrac12\mathsf O_{\mathrm T,\mathsf h}^{M_0}
        +\tfrac12\mathsf O_{\mathrm T,\mathsf h}^{M_1}\right)
 &\le\delta+16\mathrm T+31\mathrm T\mathsf h.
 \label{eq:phil-frozen-path-mixture}
\end{align}
The frozen arm keeps its original uniform rows.

For resolution \(\mathsf a>0\), let
\(\mathscr E_M(\mathsf a,\mathsf h)\) require throughout the interval
that the assignment is \(M\), the stacked residual norm
\(\sqrt{(r^b)^2+(r^m)^2}\) on its two edges is at most
\(\mathsf a\), and that norm on each exterior edge exceeds
\(\mathsf a\).
Define
\[
 \mathsf b(\mathrm T,\mathsf a,\mathsf h)
 =8\sqrt{2/\pi}\mathsf a+
            (8+3400\mathsf h)\mathrm T/\mathsf a^2.
\]
Then
\begin{equation}
 \Pr_{\mathrm L}(\mathscr E_M^c)
 \le(\mathsf h+2)\delta+\mathsf b(\mathrm T,\mathsf a,\mathsf h).
 \label{eq:phil-selected-coherence-persistence}
\end{equation}
In particular, the choice
\(\mathsf a=\mathrm T^{1/4}\),
\(\mathsf h=\mathrm T^{-1/4}\), and \(\delta\le\mathrm T\)
makes this probability tend to zero.
The matching-window label, equal to \(M_j\) when
\(\mathscr E_{M_j}\) holds and to an additional symbol otherwise,
converges to \(\delta_M\) in the learned arm and to the equal matching
mixture in the frozen arm.

At one time, define \(\mathfrak P_{\mathsf a}(\mathbf x)\) as the
connected-component partition of edges whose stacked residual norm
\(\sqrt{(r^b)^2+(r^m)^2}\) is at most \(\mathsf a\).
For times when both
one-time hypotheses above hold, put
\[
 \mathsf e(\mathrm T,\mathsf a)
 =4\sqrt{2/\pi}\mathsf a+4\mathrm T/\mathsf a^2.
\]
The learned and frozen partition laws satisfy
\begin{equation}
 \left|\operatorname{TV}
   (q_t^{\mathrm L}\mathfrak P_{\mathsf a},
                 q_t^0\mathfrak P_{\mathsf a})-\tfrac12\right|
 \le2\delta+16\mathrm T+2\mathsf e(\mathrm T,\mathsf a).
 \label{eq:phil-coherent-partition-contrast}
\end{equation}
Thus the contrast persists for the coherent partition, not only the
source label.

These conclusions apply along \cref{cor:phil-joint-selection} with
tolerance \(\delta\le\mathrm T\), mobility cap
\(\epsilon_{\mathrm L}\le\mathrm T^2\), and \(t_0=t_{\mathrm L}\).
The matching can vary along that schedule. The internal restoring rates
are \(8+\mathrm T\) and \(4+5\mathrm T\), on an event whose probability
tends to one. The resolution therefore has restoring ratio at least
\(4\mathsf a^2/\mathrm T\to\infty\).
If the rows are also within \(\delta\) of their matching vertex,
the stationary law with those rows frozen realizes the same persistence
event with probability at least
\(1-4(\mathsf h+2)\delta-\mathsf b\).
This late coherence is restored equilibrium coherence; it supplies no
entrance or first-formation assertion.
\status{ESTABLISHED}

\noindent\emph{Proof.}
First retain the proposal mechanism in the comparison. Put
\(r(t)=q_t^{\mathrm L}(S\ne M)\), and let
\(b_{\ne M}(t,\mathbf x)\) be the probability of proposing a different
assignment at a clock ring. Since the clock has rate one,
\[
 r'(t)=\E b_{\ne M}(t,\mathbf X_t)-r(t).
\]
If \(N_{\ne M}\) counts proposals to an assignment other than \(M\)
on the interval, its compensator gives
\begin{equation}
 \E N_{\ne M}
 =r(t_0+\mathsf h)-r(t_0)+
                   \int_{t_0}^{t_0+\mathsf h}r(t)dt.
 \label{eq:phil-selected-proposal-occupation}
\end{equation}
Consequently
\(\Pr(\exists t\in[t_0,t_0+\mathsf h]:S_t\ne M)
 \le(\mathsf h+1)\delta\).
Couple the initial law maximally with \(\nu_{\mathrm T}^M\), and use
the same future Brownian motion for the native process and the
stationary matching process. Their paths agree until an initial
mismatch or a proposal to another assignment. The union bound and
\cref{eq:phil-selected-proposal-occupation} prove
\cref{eq:phil-selected-path-coupling}. This argument counts proposals
of the native process and does not condition its diffusion on an
assignment history.

For the frozen comparison, use the normalized weights
\(1,\mathsf r_E,\mathsf r_N,\mathsf r_C\) from
\cref{sec:phil-pairing-selection-proof}. The determinant bounds in the
proof of \cref{prop:phil-temperature-tracking} give
\[
 \mathsf r_E,\mathsf r_N\le\tfrac52\mathrm T,
 \qquad \mathsf r_C\le\tfrac32\mathrm T^{3/2}.
\]
The invariant mass \(\kappa_{\mathrm T}\) of nonmatching assignments
is therefore at most \(16\mathrm T\). Conditional on either matching,
its configuration law is the Gaussian factor of \(\nu_{\mathrm T}^M\).
The invariant law is consequently within \(\kappa_{\mathrm T}\) of
the equal mixture of these two laws.

Detailed balance bounds the stationary flux from \(M\) to all
nonmatching assignments by \(\kappa_{\mathrm T}\).
For the other matching \(M'\), write \(\mu_M(\mathbf x)\) for the
joint invariant density on the matching sheet. Pointwise,
\[
 \mu_M b_{M'}
 =\frac{\mu_M\mu_{M'}}{\sum_S\mu_S}
 \le\tfrac12\sqrt{\mu_M\mu_{M'}}.
\]
Since \((\mathsf A_M+\mathsf A_{M'})/2\) is the oriented-cycle
matrix, Gaussian integration gives
\(\int\sqrt{\mu_M\mu_{M'}}=\mu_{\mathrm T,1/2}(M)\mathsf r_C\).
Dividing the two flux bounds by the matching mass yields
\[
 \E_{\nu_{\mathrm T}^M}b_{\ne M}
 \le8\mathsf r_E+4\mathsf r_N+\tfrac52\mathsf r_C
 \le31\mathrm T.
\]
Thus the stationary flux out of the two matching sheets together is
at most \(31\mathrm T\). Couple the stationary native process at
its initial time with the equal matching mixture and use common
Brownian increments until the first departure from its coupled
matching. Its failure probability is at most
\(16\mathrm T+31\mathrm T\mathsf h\). The frozen native process
started at \(q_{t_0}^0\) differs from its stationary path law by at
most its initial total variation \(\delta\), because the frozen
transition kernel is the same. This proves
\cref{eq:phil-frozen-path-mixture}.

It remains to bound residuals on the comparison path.
The four normal coordinates
\(z_g^\pm=n_g^\pm\cdot\mathbf X\), \(g\in M\), from
\cref{eq:rg-formation-cycle-normals} have rates
\(8+\mathrm T\) and \(4+5\mathrm T\). Their stationary variances
are at most \(\mathrm T/4\), and
\[
 r_g^b=\sqrt2z_g^++z_g^-,\qquad
 r_g^m=\sqrt2z_g^+-z_g^-.
\]
For a scalar Ornstein--Uhlenbeck noise starting at zero, integration
by parts in its stochastic convolution gives
\[
 Z_t=\sqrt{2\mathrm T}
 \left[W_t-\lambda\int_0^te^{-\lambda(t-u)}W_u du\right],
 \qquad
 \E\sup_{t\le\mathsf h}|Z_t|^2\le32\mathrm T\mathsf h.
\]
The last inequality uses
\(\sup|Z_t|\le2\sqrt{2\mathrm T}\sup|W_t|\) and the Brownian
maximal estimate proved in the matching-episode argument of
\cref{sec:rg-formation-cycle-proof}.
For the four-dimensional stationary normal vector \(z\), it follows
that
\[
 \E\sup_{t\le\mathsf h}|z_t|^2
 \le2\mathrm T+256\mathrm T\mathsf h.
\]
Each selected stacked residual norm is at most \(2|z_t|\), since
\((r_g^b)^2+(r_g^m)^2=4(z_g^+)^2+2(z_g^-)^2\).
Thus the probability that any internal stacked norm exceeds
\(\mathsf a\) at any time is at most
\((8+1024\mathsf h)\mathrm T/\mathsf a^2\).

Every exterior link coordinate is absent from \(H_M\), and its
block in \(\mathsf A_M\) is the scalar \(\mathrm T\).
It is therefore an independent \(\mathcal N(0,1)\) coordinate in
the stationary Gaussian. Each exterior belief or model residual
has variance at least one. The probability that any of these four
residuals initially has absolute value at most \(2\mathsf a\)
is at most \(8\sqrt{2/\pi}\mathsf a\).

To control their increments, diagonalize the symmetric positive
matrix \(\mathsf A_M\). In a scalar eigencoordinate with rate
\(\lambda>0\), stationarity and the preceding noise estimate imply
\[
 \E\sup_{t\le\mathsf h}|X_t-X_0|^2
 \le\frac{2\mathrm T}{\lambda}(1-e^{-\lambda\mathsf h})^2
                       +64\mathrm T\mathsf h
 \le66\mathrm T\mathsf h,
\]
since \((1-e^{-v})^2\le v\) for \(v\ge0\).
Summing twelve coordinates gives
\begin{equation}
 \E\sup_{t\le\mathsf h}|\mathbf X_t-\mathbf X_0|^2
 \le792\mathrm T\mathsf h.
 \label{eq:phil-selected-ou-increments}
\end{equation}
Every residual covector has norm \(\sqrt3\). Except with
probability \(2376\mathrm T\mathsf h/\mathsf a^2\), all exterior
increments are therefore at most \(\mathsf a\) in absolute value.
The initial small-ball estimate and the two maximal estimates prove
\(\mathsf O_{\mathrm T,\mathsf h}^M(\mathscr E_M^c)\le\mathsf b\).
Together with \cref{eq:phil-selected-path-coupling}, this proves
\cref{eq:phil-selected-coherence-persistence}.
The frozen path comparison proves the stated matching-window label
limit whenever \(\mathsf b\to0\),
\(\mathsf h\delta\to0\), and \(\mathrm T\mathsf h\to0\), in
particular for the displayed choices.

For a single time no increment estimate is needed. The normal second
moment is at most \(\mathrm T\), so the internal failure probability
is at most \(4\mathrm T/\mathsf a^2\). The exterior small-ball
bound at radius \(\mathsf a\) is
\(4\sqrt{2/\pi}\mathsf a\). Thus the partition induced by
\(\nu_{\mathrm T}^{M_j}\) differs from the matching partition
\(\mathfrak P_j\) with probability at most \(\mathsf e\).
The learned partition law is within \(\delta+\mathsf e\) of
\(\delta_{\mathfrak P_M}\); the frozen one is within
\(\delta+16\mathrm T+\mathsf e\) of
\((\delta_{\mathfrak P_0}+\delta_{\mathfrak P_1})/2\).
These matching partitions are distinct, so the triangle inequality
for total variation proves \cref{eq:phil-coherent-partition-contrast}.

The schedule in \cref{cor:phil-joint-selection} permits tolerance
\(\mathrm T\) and mobility cap \(\mathrm T^2\).
Then \(\mathrm T^{-1/4}\le t_{\mathrm L}\) for sufficiently small
\(\mathrm T\), so the following persistence window lies within the
selection window. All errors just bounded tend to zero.
The normal rates establish the restoring ratio on the native event.
Finally, if \(\max_i|p_i-p_i^M|\le\delta\),
\(\operatorname{TV}(\mu_{\mathrm T,p},\nu_{\mathrm T}^M)
 =1-\mu_{\mathrm T,p}(M)\le4\delta\).
Applying the same occupation and coupling argument to this
stationary frozen-row process gives the stationary lower bound in
the statement. The persistence event is therefore already typical
at equilibrium for those rows and carries no first-entrance claim.
\hfill\(\square\)


